\section{Related Work}

\subsection{Deep Learning for Skin Lesion Classification}

The application of deep learning to dermoscopic image analysis gained
significant attention following Esteva et al.'s demonstration that a CNN
trained on clinical images could match dermatologist-level performance
\cite{esteva2017dermatologist}. Subsequent work has explored various
architectures and training strategies. The ISIC challenges have driven
progress in this area by providing standardized benchmarks and large-scale
annotated datasets. Recent approaches have employed transfer learning from
ImageNet-pretrained models, with EfficientNet variants achieving strong
results due to their favorable accuracy-efficiency tradeoff
\cite{tan2019efficientnet}. Gessert et al. \cite{gessert2020skin} demonstrated
that ensembles of EfficientNet models with extensive data augmentation could
achieve top performance on ISIC challenges, though their work focused
primarily on classification accuracy rather than interpretability.

\subsection{Explainability in Medical Imaging}

Explainable AI (XAI) has become increasingly important in medical applications
where understanding model reasoning is critical for clinical acceptance.
Gradient-weighted Class Activation Mapping (Grad-CAM) \cite{selvaraju2017grad}
and its extensions provide visual explanations by highlighting image regions
that contribute most to predictions. In dermatology, attention visualization
helps verify that models focus on lesion features rather than artifacts or
irrelevant background regions. Chattopadhay et al. \cite{chattopadhay2018grad}
proposed Grad-CAM++, which provides more accurate localization through a
weighted combination of positive partial derivatives, making it particularly
suitable for medical imaging where precise localization matters.

Concept-based explanations offer an alternative approach by relating model
predictions to human-understandable concepts. Kim et al.
\cite{kim2018interpretability} introduced Testing with Concept Activation
Vectors (TCAV), which learns directions in a network's feature space
corresponding to user-defined concepts. This approach has been applied to
medical imaging to explain predictions in terms of clinically meaningful
attributes, and recent work such as FastCAV \cite{schmalwasser2025fastcav} has
reduced its computational cost. Our SGD-CAV implementation adapts this framework
for efficient concept-based explanations aligned with dermatological concepts.

\subsection{Uncertainty Quantification in Deep Learning}

Reliable uncertainty estimates are essential for clinical decision support
systems, as they enable appropriate human-AI collaboration by identifying
cases requiring expert review. Gal and Ghahramani \cite{gal2016dropout}
demonstrated that dropout applied at test time approximates Bayesian
inference, providing uncertainty estimates without architectural changes.
This Monte Carlo Dropout approach has been widely adopted in medical imaging
due to its simplicity and effectiveness. Subsequent work has decomposed
uncertainty into epistemic (model uncertainty) and aleatoric (data
uncertainty) components \cite{kendall2017uncertainties}, providing more
nuanced characterization of prediction confidence. Our system implements
this decomposition to distinguish between cases where the model lacks
knowledge versus cases with inherently ambiguous features.
